\section{The record, and what each document contains}\label{sec:corpus}

\subsection{Five editions under one concept DOI}\label{ssec:editions}

The record 10.5281/zenodo.22883643 is the latest of five versions of the concept record 10.5281/zenodo.22666540. Each version kept the files of the one before and added new ones, so the latest version holds all 43 files. The version history, read from the record's public metadata on 26 September 2026:

\begin{center}\small
\begin{tabularx}{\linewidth}{@{}>{\raggedright\arraybackslash}p{0.2\linewidth}>{\raggedright\arraybackslash}p{0.08\linewidth}>{\raggedright\arraybackslash}X@{}}
\toprule
Version (DOI 10.5281/zenodo.\dots) & Files added & What was added\\
\midrule
22666541, 8 Sept.\ (\emph{2026.09.08}) & 6 & The first checkpoint: a 46-page quantum manuscript (\fn{01\_}), a 72-page spatial manuscript (\fn{02\_}), their sources (\fn{03\_}), a sealed author packet (\fn{04\_}), the record's \fn{README.md} and a manifest\\
22667605, 8 Sept.\ (\emph{full-workbench}) & 7 & Later snapshots of the quantum and spatial lines (\fn{05\_}, 53 pp.; \fn{06\_}, 108 pp.) and the first volume manuscript (\fn{07\_}, 62 pp.); the whole Overleaf workbench (\fn{08\_}, 87~MB); a reading map, an edition guide and a manifest (\fn{09\_}--\fn{11\_})\\
22678364, 9 Sept.\ (\emph{consolidated}) & 10 & The five foundation readers (§\ref{ssec:lines}); the sources of 9 September; an index for machine reading; manifests; a reading guide\\
22803564, 21 Sept.\ (\emph{heat-volume}) & 14 & The 297-page reader of the 14--16 September continuation; the 49-page reader of 17 September; the 46-page heat and volume reader of 21 September with its sources; two cumulative archives; reading guides; the attribution clarification\\
22883643, 21 Sept.\ (\emph{fifth-reference}) & 6 & The 24-page fifth-reference reader, its LaTeX, Markdown and sources, a cumulative archive, and its reading guide\\
\bottomrule
\end{tabularx}
\end{center}

Two facts matter for a visitor.
\begin{itemize}
\item The record's front file \fn{README.md} was written for the first version, and it describes the 8 September checkpoint. The current guide is \fn{YM\_READING\_GUIDE\_20260921\_FIFTH.md}, and the current reader is \fn{yang\_mills\_fifth\_reference\_20260921.pdf}.
\item Twenty-three of the files are byte-identical to files in the repository at commit \fn{fa79faf} (checked by MD5, Appendix~\ref{app:verification}). The other twenty are not files of the repository, but five of them occur, byte-identical, inside archives committed there: \fn{AI\_READING\_INDEX.md}, \fn{READER\_PROVENANCE.json} and \fn{SOURCE\_MANIFEST.json} in \fn{yang\_mills\_current\_sources\_2026-09-09.zip}, and \fn{ALPOGE\_ATTRIBUTION\_20260921.md} and \fn{YM\_READING\_GUIDE\_20260921.md} in \fn{yang\_mills\_heat\_volume\_reader\_sources\_20260921.zip}.
\end{itemize}

\subsection{The documents, line by line}\label{ssec:lines}

The mathematical documents form three lines of manuscripts from 8--9 September, each with several dated snapshots, followed by four readers of later continuations. In each line the section list of the latest snapshot contains those of the earlier ones, and the 9 September guide says that the quantum reader is cumulative.

\begin{center}\small
\begin{longtable}{@{}>{\raggedright\arraybackslash}p{0.25\linewidth}>{\raggedright\arraybackslash}p{0.2\linewidth}>{\raggedright\arraybackslash}p{0.49\linewidth}@{}}
\toprule
Document & Files in the record & Content, in the document's own terms\\
\midrule
\endhead
\emph{Quantum line.} \emph{Extensive true-vacuum quantum blocking in the full spatial Wilson box} (8--9 Sept.) & \fn{01\_} (46 pp.), \fn{05\_} (53 pp.), \fn{quantum\_nonabelian\_vertex} (63 pp.), \fn{quantum\_interacting\_tensor\_band} (71 pp.), \textbf{\fn{quantum\_coarse\_graining}} (79 pp., cumulative) & Exact blocking maps of the finite interacting Hamiltonian; a uniform single-link estimate for the actual vacuum; zero-shift reconstruction and comparison of low-energy states; the scope of the spectral conclusion (§11); the material-tensor transfer of the Jacobian map into local-energy states (§§13--16, with Theorems 13.1 and 14.2); a nonzero non-Abelian vertex (§17); the Navier--Stokes curvature-state map (§§20--22)\\
\emph{Spatial line.} \emph{The retained cusp, spatial refinement, and the spectral measure of the actual magnetic vacuum translation} (8--9 Sept.) & \fn{02\_} (72 pp.), \fn{06\_} (108 pp.), \textbf{\fn{spatial\_continuum}} (129 pp.) & Joint geometric and coupling sequences built from the $S^6$ period data; spectral escape and its failure modes; weak-coupling variational bounds; exact plaquette-to-loop maps and a pure-gauge configuration limit; covariance and magnetic-mass lower bounds; a fluid-to-curvature map; §27, the collapse of the finite-box weak-coupling disk (Proposition~\ref{prop:neumann}); §28, the current conclusion\\
\emph{Volume line.} \emph{Volume-independent local estimates in the full SU(2) Wilson vacuum} (8--9 Sept.) & \fn{07\_} (62 pp.), \textbf{\fn{volume\_uniform\_vacuum}} (81 pp.) & Local projections and comparison operators; local electric moments at every positive coupling; uniform local Wilson and electric covariances; cluster coordinates; the gauge-restricted free threshold $3\kappa$; negative-period collapse; the product remainder on the cusp at $0<\xi\le10^{-16}$; the logarithmic coupling path\\
Web continuation, 14--16 Sept. & \fn{yang\_mills\_web\_continuation} (297 pp.) and its sources & The source equations, gauge-coordinate maps, cubic calculation and spectral return; the first gap theorems (order 1 and order 2; Section~\ref{sec:earlier})\\
Quartic cube reader, 17 Sept. & \fn{yang\_mills\_quartic\_cube\_reader} (49 pp.) and its sources & The fourth-order vacuum source in two presentations, equal as polynomials; the sixth-order energy with the cube contribution $-83/1944$; a gap bound for $g^2\ge15/4$\\
Heat and volume reader, 21 Sept. & \fn{yang\_mills\_heat\_volume\_20260921} (46 pp.) and archives & Fourth-order heat correlations; a box-independent heat remainder on the circle $|\xi|<3/256$; a fixed-spacing spatial limit; the plaquette family's coupling to the whole complementary space\\
\textbf{Fifth-reference reader}, 21 Sept. & \textbf{\fn{yang\_mills\_fifth\_reference\_20260921}} (24 pp.) and archive & The latest results: the gap bound of Section~\ref{sec:gap} and the heat and complement bounds of Section~\ref{sec:heat}, with both proof manuscripts (F1--F42, H1--H34)\\
\bottomrule
\end{longtable}
\end{center}
Bold marks the current document of each line. The cumulative archives of 21 September also hold the continuation of 18--19 September (sixth source, eighth energy, a 240-face response; directory \fn{20260918-sixth-source}); the workbench says that its new spatial tables and final execution package were not delivered (repository \fn{README.md}).

\subsection{How the documents describe their own scope}\label{ssec:scope}

\begin{itemize}
\item The 9 September readers are finite-regulator calculations. The spatial reader's conclusion (§28) says that the actual interacting continuum construction, and nonzero low-energy spectral weight outside its proved exclusions, remain unresolved, and that no missing covariance, tightness or convergence statement was assumed.
\item The quantum reader's §11 records that rewriting a coarse kinetic coefficient at a \emph{running coupling} is only a coefficient dictionary, not a renormalisation of the original Hamiltonian.
\item The volume reader's cusp estimates hold at fixed coupling $0<\xi\le10^{-16}$, that is $g^2\ge5\cdot10^7$; its §24 treats the logarithmic coupling path without a small-$\xi$ gap.
\item Every reader of 17 and 21 September states that no four-dimensional continuum mass gap is claimed.
\end{itemize}
These statements were checked against the documents where this paper cites them. The results of the 8--9 September lines were not otherwise audited here; Section~\ref{sec:earlier} lists the few that were.
